\documentclass[11pt]{article}
\usepackage[margin=1in]{geometry}
\usepackage{amsmath,amssymb,booktabs}
\title{Canonical Nondimensionalization of the Kelvin--Voigt Dispersion Relation}
\author{}
\date{}
\begin{document}
\maketitle

\section{Dimensional relation}
Let $\eta_i=\mu_i+G_i/\gamma$ and $\rho_T=\rho_1+\rho_2$. The independently
transcribed starting equation is
\begin{equation}
\gamma\left[
\frac1{\eta_1+\eta_2\sqrt{1+\rho_2\gamma/(\eta_2k^2)}}+
\frac1{\eta_2+\eta_1\sqrt{1+\rho_1\gamma/(\eta_1k^2)}}
\right]+\frac{4k^2}{\rho_T}=0.\label{eq:dim}
\end{equation}
Expanding $\eta_i$ gives the supplied equation exactly.

The dimensions are $[\gamma]=T^{-1}$, $[k]=L^{-1}$,
$[\rho_i]=ML^{-3}$, $[\mu_i]=ML^{-1}T^{-1}$, and
$[G_i]=ML^{-1}T^{-2}$. Eight quantities and a rank-three dimension matrix give
five independent $\Pi$ groups.

\section{General clock}
Set
\[
\gamma=s/t_c,\quad \rho_i=\rho_Tr_i,\quad \mu_i=\mu_Tm_i,
\quad G_i=G_Tg_i,
\]
and define
\[
p=\frac{\rho_T}{\mu_Tk^2t_c}=\frac{t_v}{t_c},\qquad
q=\frac{G_Tt_c}{\mu_T},\qquad H_i=m_i+\frac{qg_i}{s}.
\]
Then
\[
\eta_i=\mu_i+\frac{G_i}{\gamma}
=\mu_Tm_i+\frac{G_Tg_it_c}{s}=\mu_TH_i
\]
and
\[
1+\frac{\rho_i\gamma}{\eta_i k^2}
=1+\frac{\rho_Tr_i(s/t_c)}{\mu_TH_i k^2}
=1+\frac{pr_is}{H_i}.
\]
Multiplying Eq.~\eqref{eq:dim} by $\rho_T/k^2$ therefore gives
\begin{equation}
ps\left[
\frac1{H_1+H_2\sqrt{1+pr_2s/H_2}}+
\frac1{H_2+H_1\sqrt{1+pr_1s/H_1}}
\right]+4=0.\label{eq:master}
\end{equation}
Moreover
\[
pq=\frac{\rho_TG_T}{\mu_T^2k^2}\equiv\Lambda.
\]
With $a_i=sH_i=m_is+qg_i$, Eq.~\eqref{eq:master} is, for $s\ne0$,
\begin{equation}
ps^2\left[
\frac1{a_1+a_2\sqrt{1+pr_2s^2/a_2}}+
\frac1{a_2+a_1\sqrt{1+pr_1s^2/a_1}}
\right]+4=0.\label{eq:cleared}
\end{equation}

\section{Clocks}
Define
\[
t_v=\frac{\rho_T}{\mu_Tk^2},\qquad
t_e=\frac1k\sqrt{\frac{\rho_T}{G_T}},\qquad
t_R=\frac{\mu_T}{G_T},\qquad
\chi=\sqrt\Lambda=\frac{t_v}{t_e}.
\]
They satisfy $t_R=t_e^2/t_v$. Their substitutions into
Eq.~\eqref{eq:cleared} are
\begin{center}
\begin{tabular}{cccc}
\toprule
$t_c$ & response & $p$ & $q$\\
\midrule
$t_v$ & $s_v$ & $1$ & $\Lambda$\\
$t_e$ & $s_e$ & $\chi$ & $\chi$\\
$t_R$ & $s_R$ & $\Lambda$ & $1$\\
$1/(r_v+r_e)$ & $s_+$ & $1+\chi$ & $\chi^2/(1+\chi)$\\
\bottomrule
\end{tabular}
\end{center}
Thus, for example, the recommended viscous form is
\begin{equation}
s_v^2\left[
\frac1{a_1+a_2\sqrt{1+r_2s_v^2/a_2}}+
\frac1{a_2+a_1\sqrt{1+r_1s_v^2/a_1}}
\right]+4=0,
\quad a_i=m_is_v+\Lambda g_i.\label{eq:viscous}
\end{equation}
The elastic and relaxation equations follow from the second and third rows
without any further assumptions.

For $E_k=(r_v-r_e)/(r_v+r_e)=(1-\chi)/(1+\chi)$,
\[
\chi=\frac{1-E_k}{1+E_k},\quad
p=\frac2{1+E_k},\quad q=\frac{(1-E_k)^2}{2(1+E_k)}.
\]
With $p_+=1+\chi$, $q_+=\chi^2/(1+\chi)$, and
$a_i=m_is_++q_+g_i$, the rate-sum form is explicitly
\[
p_+s_+^2\left[
\frac1{a_1+a_2\sqrt{1+p_+r_2s_+^2/a_2}}+
\frac1{a_2+a_1\sqrt{1+p_+r_1s_+^2/a_1}}
\right]+4=0.
\]

\section{Symmetric coordinates and recommendation}
Use
\[
r_{1,2}=\frac{1\pm A_\rho}{2},\qquad
m_{1,2}=\frac{1\pm A_\mu}{2},\qquad
g_{1,2}=\frac{1\pm A_G}{2}.
\]
Together, $(s_v,\Lambda,A_\rho,A_\mu,A_G)$ are five independent
dimensionless coordinates. The three contrasts and $E_k$ are convenient
bounded transforms of ratio-based groups, not a unique monomial Buckingham
basis. Equation~\eqref{eq:viscous} is recommended because it is symmetric,
minimal, algebraically simplest, and remains defined as $G_T\to0$.

\end{document}
